%% ===================================================================
\section{Introduction}
\label{sec:intro}

Network planning with a digital twin requires a propagation model for places
where no measurement has been taken yet: coverage, interference and site
selection must be assessed before a campaign exists in the target
region~\cite{khan2022dt,alkhateeb2023dt}. Ray tracing resolves the geometry
but must be rerun for every change of layout and needs detailed 3-D
data~\cite{yun2015raytracing,hoydis2023sionna}; empirical formulas such as
Okumura--Hata or 3GPP TR~38.901 are cheap but apply one parameter set to a
whole environment class~\cite{hata1980,tr38901}. Tile-based models sit
between the two: the twin is partitioned into tiles, each tile carries the
two parameters $(\gamma,P_{L_0})$ of a log-distance law, and parameters
learned where measurements exist are transferred to unmeasured tiles by
environmental similarity~\cite{ozyurt2020maple,ozyurt2026}. The transported
quantity is small and physically interpretable, and a planning tool can
evaluate it for any site layout.

Whether such transfer works \emph{across cities}, and what its errors cost a
network planner, has not been established. The descriptor-based framework
of~\cite{ozyurt2026}, which transfers tile parameters with a Gaussian kernel
over a descriptor space, was evaluated on two urban datasets that are not
public. Three questions remain open. First, how reliable are the tile labels
themselves, given that a tile of \SI{100}{\metre} sees its serving sites over
a narrow range of distances and that neighbouring receivers share
propagation paths? Second, how does kernel transfer behave when the target
city differs from every source city, which is exactly when transfer is
needed, and how does it compare with learned regressors under one protocol?
Third, what do parameter errors mean for the quantities a planner uses:
path loss, serving cell, SINR, coverage, cell edges and the choice of sites,
is the tile parameter the right quantity to transfer at all, and how much
coverage can an operator promise from a transferred model?

This paper answers these questions on a public benchmark. Our contributions
are the following.
\begin{enumerate}
\item \textbf{Reproducible multi-city, multi-band corpora.} Eleven
geo-referenced city scenes built from a pinned release of open building data
are ray traced from \nSites{} rooftop sites at \SI{\freqGHz}{\giga\hertz}
(\nSamples{} links, \nTiles{} labelled tiles) and at
\SI{\freqTwoGHz}{\giga\hertz} in the upper mid-band, with a ray-budget study,
a physical consistency check of the solver output, a replicate with random
rooftop siting and an external test on independently modelled scenes. Every
number in this paper is regenerated from the released data by one command.
\item \textbf{Tile labels with honest precision.} We show that the per-tile
fit is ill-conditioned (Proposition~\ref{prop:identify}), replace ad-hoc
regularisation by empirical-Bayes shrinkage whose risk is bounded whatever
the tile geometry (Proposition~\ref{prop:eb}) with a small-sample (CR2)
cluster-robust variance, and show that the labels depend on the serving
sites, the tile size and the dynamic range of the data.
\item \textbf{An analysis of kernel transfer.} We prove that the Gaussian
kernel operator has two limiting regimes (Theorem~\ref{thm:regimes}),
measure where the tuned operator lies on every held-out city, and show that
learning its metric removes most of its disadvantage.
\item \textbf{A leave-one-city-out benchmark with city-level inference.}
Standard-model, label-averaging, Gaussian-process, linear, tree, neural and
domain-adaptation operators are compared under one tuning protocol with five
seeds, with and without surveyed tiles (random, contiguous or designed), and
every claim is tested at the level of cities with bootstrap intervals,
Wilcoxon tests and Holm correction.
\item \textbf{Network-level evaluation and site-aware link transfer.} We
derive what a tile model can express for cell association and SINR
(Proposition~\ref{prop:tilemodel}), evaluate every operator through path
loss, serving cell, SINR, coverage, cell edges, dimensioning and site
planning scored on the ray-traced truth, and propose a hurdle link model
conditioned on the obstruction geometry of each link that outperforms even
the tile model with perfect labels in every held-out city.
\item \textbf{Survey design and coverage certificates.} We give a survey
design for MDT-style site calibration with a $(1-1/e)$ guarantee
(Proposition~\ref{prop:survey}) and a distribution-free, city-level coverage
certificate from leave-one-city-out residuals
(Proposition~\ref{prop:certify}), and measure what each transfer model can
promise at a stated risk.
\end{enumerate}

Section~\ref{sec:related} reviews related work, Section~\ref{sec:model}
formalises the network model and the transfer problem,
Section~\ref{sec:methods} describes the corpus, labels, descriptors,
operators and evaluation, Section~\ref{sec:results} reports the results,
Section~\ref{sec:discussion} discusses their implications and limitations,
and Section~\ref{sec:conclusion} concludes. Proofs are in \ref{app:proofs}
and the complete city-level comparisons in \ref{app:stats}.
